\section{Complete implementation and measured simplification}
\label{sec:implementation}

\subsection{Input contract and trusted code}
The primary program reads a JSON object with exactly two fields,
\texttt{goal} and \texttt{proof}. The goal is a closed canonical formula.
The checker prints \texttt{ACCEPT} only after inferring the proof's
conclusion, validating its syntax and scope, and comparing it with the
independently validated goal. An exception or nonzero process exit rejects.
No Python \texttt{assert} is used in the kernel, so optimized execution
does not disable a logical check.

Zero is integer \texttt{0}; negative integers are free parameter levels;
positive integers are bound indices, with index one innermost. Terms
\texttt{[0,t]} and \texttt{[n,a,h]} encode successor and iteration, with
one bound index added only inside \texttt{h}. A two-element formula
\texttt{[s,t]} is an equation; \texttt{[B]} binds two indices in its
matrix, with the counter at index one and the witness at index two.
The sorted syntax traversal distinguishes a successor from an equation.
Booleans, floats, out-of-scope indices, and noncanonical terms are rejected.

The instruction formats are
\[
\begin{array}{ll}
 \texttt{["u",j,t,p]}&\text{use clause }j\text{ at }t,\\
 \texttt{["i",B,t,p]}&\text{introduce }[B]\text{ at witness }t,\\
 \texttt{["e",p,q]}&\text{infer a major from }p\text{ and check branch }q.
\end{array}
\]
For unconditional clause use, the fourth payload is unused; the examples
use JSON \texttt{null}. This is a declared certificate convention, not
an unvalidated proof premise. Conditional use must infer exactly the
retained antecedent. The context initially contains only
\texttt{([], [1,1])}, the reflexivity clause. Every I or O$^*$ branch adds
one clause and one parameter, so the context length determines the next
fresh level without a separate numerical-scope argument.

Every successful recursive call validates its inferred formula in that
call's outer scope. Thus opening cannot export its fresh parameter,
including when the synthesized term contains it. Introduction returns
the unchanged annotated matrix and validates that matrix at the end;
malformed syntax temporarily erased in a computed branch cannot survive
this final check. The independent validation includes such erased
malformation examples.

\paragraph{The retained-antecedent invariant.}
The implementation computes the original antecedent $u=N(B(b,0))$ and
substitutes the synthesized name into the body. It need not recompute
the antecedent. If the probe succeeds, $u$ is literally $H$, and
substitution leaves it unchanged. If it fails, the substitution is the
identity. Therefore the stored pair is exactly the clause stipulated
by O$^*$ in both cases. This identity accounts for part of the matched
opening-control comparison below.

\subsection{Source measurements}
The measurement convention follows the repository: remove a module's
documentation string, parse the remaining source, normalize it with
CPython 3.12.14 \texttt{ast.unparse}, and append one newline. Count UTF-8
bytes, significant Python tokens (excluding comments, indentation,
newlines and end markers), and all nodes from \texttt{ast.walk}.
The gzip column uses level 9, a zero timestamp, and these normalized
bytes. Gzip is an informational metric, not an optimization objective.
All trusted syntax, normalization, binding, scope, proof, and JSON
handling code is included. No untrusted proof-construction code is
required for acceptance.

\begin{table}[htbp]
\centering\small
\begin{tabular}{@{}lrrrr@{}}
\toprule
Complete program & Bytes & Tokens & AST & Gzip 9\\
\midrule
Inferred opening (selected) & 1,990 & 683 & 771 & 819\\
Matched equality opening & 1,998 & 676 & 765 & 812\\
Matched equality output & 1,983 & 663 & 751 & 812\\
Incoming witness discharge & 2,273 & 732 & 813 & 903\\
Incoming uniform binding & 2,238 & 693 & 764 & 866\\
Incoming clause contexts & 2,282 & 666 & 770 & 774\\
Inferred opening with fuel & 2,203 & 728 & 826 & 885\\
Incoming witness with fuel & 2,486 & 777 & 868 & 967\\
\bottomrule
\end{tabular}
\caption{Complete programs under one reproducible source-measurement convention.}
\label{tab:metrics}
\end{table}


Against the incoming explicit-equality discharge implementation, the
selected program saves 283 normalized bytes (12.45\%), 49 significant
tokens (6.69\%), and 42 AST nodes (5.17\%); gzip falls by 84 bytes.
This comparison includes both the new primitive interface and the
change to a forward checker that infers conclusions once. It must not
be attributed entirely to the mathematical rule restriction.

The matched explicit-equality \emph{opening} control uses the same
syntax traversal, normalizer, context convention, and forward interface.
Its opening accepts an equality proof $b=t$, specializes the clause at
$t$, and checks the branch. The new checker saves eight bytes, while
using seven more significant tokens, six more AST nodes, and seven more
gzip bytes. The matched explicit-equality \emph{output} control instead
keeps the unspecialized clause and substitutes into the branch's result;
it is seven normalized bytes shorter than the selected checker.
These mixed results are recorded rather than concealed: removing a
proof premise does not automatically shorten all implementations.

The selected checker is shorter in normalized bytes than the six
inspected explicitly typed incoming checkers, but not in every coordinate.
For example, \texttt{uniform\_binding} has seven fewer AST nodes, and
\texttt{clause\_contexts} has fewer tokens and one fewer AST node.
The repository's 516-byte compact codec and 788-byte validated-public-goal
implementation remain substantially shorter~\cite{repo}. Their
representation and internal validation contracts differ, so they are
not matched controls; their smaller sizes nonetheless preclude a
global source-size record claim for this package.

\subsection{The fuel-bounded certificate relation}
The second standalone checker adds an integer \texttt{fuel} field. It
decrements fuel on every syntax/normalization traversal entry and proof
call, rejecting before the counter becomes negative. Its full size is
included in the table. For each accepted unbounded certificate, its
finite computation uses some finite amount of fuel; conversely every
fuelled acceptance follows the same checking computation. The two
idealized checkers therefore accept proofs of exactly the same formulas
when sufficient resources are supplied.

With input size and fuel fixed, the number of recursive entries is
bounded. Each entry performs finite operations on syntax whose growth
has a primitive-recursive bound in those two inputs. The fuelled
certificate relation is consequently primitive recursive. This does
not provide a primitive-recursive bound on the fuel needed to normalize
all descriptions, and does not claim that iPRA proves the uniform
normalizer total. As with the baseline Python programs, implementation
stack, memory, and parser limits can reject a large mathematically
valid certificate; such rejection never licenses acceptance of an
invalid one.

\subsection{Executed evidence}
The final primary source has SHA-256
\begin{center}\footnotesize\ttfamily
79509b47cbbf20df923cc7b3e86952df33dcb3ce15b71f8e1d798212805a1902
\end{center}
The reproduction suite accepts 12 concrete closed certificates and
rejects 38 adversarial objects. Positive examples include both directions
of conditional symmetry, automatic transport at zero and nonzero
witnesses, nested iterator binders, nested paired formulas, a branch
that genuinely requires the specialized clause, and a synthesized term
containing the fresh parameter. The suite performs 12 exact translations
into the explicit-opening control, 24 exact-fuel/one-less checks, and
114 command-line checks with and without Python optimization. It also
checks a case in which outer substitution creates the naming guard:
blindly retaining the old branch fails, while the branch prescribed by
the substitution theorem succeeds.

An independent validator uses named bound variables, a separately
written sorted parser, and numerical evaluation by ordinary loops.
It passes 6,292 checks with no failures, including 700 normalization
comparisons, 1,398 numerical semantic comparisons, 500 binder-opening
comparisons, 400 simultaneous substitutions, three sets of 350
probe/clause/invariant checks, 1,500 arbitrary-JSON grammar comparisons,
and targeted scope and malformed-input examples. The source hash and
complete category counts are retained in the evidence JSON.

These executions address concrete implementation risks: capture,
accidental Python numeric coercion, malformed erased syntax, wrong
equality orientation, a nonzero reflexive base, a compound matrix,
self-containing substitutions, and eigenparameter escape. They are
finite tests. The universal soundness, substitution, completeness, and
bi-interpretation conclusions rest on the mathematical proofs and the
attributed inherited bridge, not on the number of successful tests.
